%%%%%%%%%%%%%%%%%%%%%%%%%%%%%%%%%%%%%%%%%%%%%%%%%%%%%%%%%%%%%%%%%%%%%%
% Lớp 9 | Chương I. Phương trình và hệ hai phương trình bậc nhất hai ẩn | Giải bài toán bằng cách lập hệ phương trình
% Dạng 9D1?3-1: Toán chuyển động
%%%%%%%%%%%%%%%%%%%%%%%%%%%%%%%%%%%%%%%%%%%%%%%%%%%%%%%%%%%%%%%%%%%%%%


% nguồn: 2026-11-14_09_de_tuyen_sinh_10_tphcm_2025_2026 (nhập 14/11/2026)
\begin{bt}%[9D1C3-1]%[ma:005006ba]
% nguon: Đề thi chính thức Tuyển sinh 10 TP.HCM năm học 2025-2026 – Bài 6
Từ vị trí $A$ của một công viên có dạng hình vuông $ABCD$ cạnh $a\text{ (km)}$, hai bạn Hoà và Bình bắt đầu chạy bộ cùng lúc với vận tốc không đổi dọc theo các cạnh của hình vuông và theo hai hướng khác nhau (Hoà chạy theo hướng $A \to B \to C \to D$, Bình chạy theo hướng $A \to D \to C \to B$). Biết rằng hai bạn gặp nhau lần thứ nhất tại vị trí $E$ trên cạnh $CD$ cách $A$ một khoảng bằng $1\text{ km}$ ($AE = 1\text{ km}$) và gặp lại nhau lần thứ hai tại vị trí $F$ trên cạnh $AB$ cách $A$ một khoảng bằng $0{,}4\text{ km}$ ($AF = 0{,}4\text{ km}$). Gọi $x, y\text{ (km/h)}$ lần lượt là vận tốc của Hoà và Bình.
\begin{enumerate}[a)]
	\item Chứng minh rằng $\dfrac{x}{y} = \dfrac{AB + BC + CE}{AD + DE}$.
	\item Tìm giá trị của $a$.
\end{enumerate}
\loigiai{
\begin{enumerate}[a)]
	\item Hai bạn xuất phát cùng lúc từ $A$. Khi gặp nhau lần thứ nhất tại $E$ trên cạnh $CD$:
	\begin{itemize}
		\item Quãng đường Hoà đã chạy là: $s_1 = AB + BC + CE$.
		\item Quãng đường Bình đã chạy là: $s_2 = AD + DE$.
	\end{itemize}
	Vì thời gian chạy của hai bạn từ lúc xuất phát đến lúc gặp nhau lần đầu là như nhau nên:
	$t_1 = \dfrac{s_1}{x} = \dfrac{s_2}{y} \Leftrightarrow \dfrac{AB + BC + CE}{x} = \dfrac{AD + DE}{y}$.\\
	Suy ra $\dfrac{x}{y} = \dfrac{AB + BC + CE}{AD + DE}$.
	\item Chu vi công viên hình vuông là $4a\text{ (km)}$.\\
	Tổng quãng đường hai bạn chạy được tính đến lần gặp thứ nhất là:
	$s_1 + s_2 = (AB + BC + CE) + (AD + DE) = AB + BC + CD + DA = 4a\text{ (km)}$.\\
	Do đó thời gian chạy đến lần gặp thứ nhất là:
	$t_1 = \dfrac{4a}{x + y}$.\\
	Từ lần gặp thứ nhất đến lần gặp thứ hai tại $F$, hai bạn tiếp tục chạy ngược chiều quanh chu vi hình vuông và gặp nhau, nên tổng quãng đường hai bạn chạy thêm đúng bằng một vòng chu vi hình vuông là $4a\text{ (km)}$.\\
	Do đó, thời gian từ lúc xuất phát đến khi gặp nhau lần thứ hai là:
	$t_2 = \dfrac{4a + 4a}{x + y} = \dfrac{8a}{x + y} = 2t_1$.\\
	Vì vận tốc của Bình là không đổi nên quãng đường Bình chạy đến lần gặp thứ hai gấp đôi quãng đường Bình chạy đến lần gặp thứ nhất.\\
	Tại lần gặp thứ hai ở $F$ trên cạnh $AB$ cách $A$ là $AF = 0{,}4\text{ km}$, quãng đường Bình đã chạy từ $A$ qua $D, C, B$ đến $F$ là:
	$s'_2 = AD + DC + CB + BF = 4a - AF = 4a - 0{,}4\text{ (km)}$.\\
	Do $s'_2 = 2s_2$, ta có:
	$4a - 0{,}4 = 2(AD + DE) = 2(a + DE) \Leftrightarrow 4a - 0{,}4 = 2a + 2DE \Leftrightarrow 2DE = 2a - 0{,}4 \Leftrightarrow DE = a - 0{,}2$.\\
	Tam giác $ADE$ vuông tại $D$ có $AD = a$, $DE = a - 0{,}2$ và cạnh huyền $AE = 1\text{ km}$. Áp dụng định lí Pythagore:
	$AD^2 + DE^2 = AE^2 \Leftrightarrow a^2 + (a - 0{,}2)^2 = 1^2 \Leftrightarrow a^2 + a^2 - 0{,}4a + 0{,}04 = 1 \Leftrightarrow 2a^2 - 0{,}4a - 0{,}96 = 0$.\\
	Chia cả hai vế cho $2$: $a^2 - 0{,}2a - 0{,}48 = 0$.\\
	Ta có $\Delta' = (-0{,}1)^2 - 1\cdot(-0{,}48) = 0{,}01 + 0{,}48 = 0{,}49 = (0{,}7)^2$.\\
	Phương trình có hai nghiệm:
	$a_1 = 0{,}1 + 0{,}7 = 0{,}8$ (thoả mãn $a > 0{,}2$),\\
	$a_2 = 0{,}1 - 0{,}7 = -0{,}6$ (loại vì $a > 0$).\\
	Vậy $a = 0{,}8\text{ (km)}$.
\end{enumerate}
}
\end{bt}
